\chapter{O que o software produz}

\eyebrow{produtos gerados}

\vspace{4pt}
Os produtos do Apex convergem para um único fim: um conjunto de dados de pesquisa
com pareamento verificável, enriquecido com medições objetivas. Este capítulo
resume os produtos gerados e o nível de confiança associado a cada um.

\section{O conjunto de dados}

O produto central é o \textbf{conjunto de dados pareado}: para cada imagem, um
vínculo verificável com sua carcaça, animal, tratamento e estrato, além da
referência física medida no animal. Ele reside no arquivo \tele{carcass.db} e é
exportado como \tele{manifest.csv} e \tele{integrity\_report.txt}.

\section{Medições por imagem}

\begin{center}
\begin{tabularx}{\linewidth}{@{}l l X@{}}
\toprule
\textbf{Medição} & \textbf{Status} & \textbf{O que é} \\
\midrule
Fração de gordura \& máscara & \pillok{validado} & Fração da superfície da carcaça
      segmentada como gordura, com sobreposição ciano. IoU $\sim$0,92. \\
Convexidade de conformação & \pillok{objetivo} & Convexidade por região (perna /
      lombo / paleta) e um mapa de convexidade. Analítica, invariante, sem
      necessidade de treinamento. \\
Grau de acabamento & \pillalert{experimental} & Uma estimativa de acabamento em
      3 classes a partir de um modelo treinado com n=22. Não validado. \\
Grau de conformação & \pillalert{estimativa} & Uma estimativa em 5 classes
      derivada da convexidade da perna. Não validado. \\
Espessura de gordura (EG) & \pillalert{experimental} & Uma estimativa por
      regressão. Não validada. \\
\bottomrule
\end{tabularx}
\end{center}

\section{Produtos por lote}

\begin{itemize}[leftmargin=1.3em,itemsep=3pt]
  \item \textbf{Estatísticas de gordura}: média, desvio padrão, mediana,
        mínimo, máximo e o histograma de distribuição.
  \item \textbf{Concordância entre avaliadores}: $\kappa$ de Fleiss por eixo e
        um grau de consenso.
  \item \textbf{Progresso da amostra}: cobertura em relação à meta de
        100--120, por estrato.
  \item \textbf{Exportação CSV}: uma linha por imagem com todas as medições.
\end{itemize}

\section{Produtos ao vivo}

\begin{itemize}[leftmargin=1.3em,itemsep=3pt]
  \item \textbf{Sobreposição de gordura ao vivo}: guia de enquadramento em
        tempo real a partir da câmera ou de vídeo (aproximada).
\end{itemize}

\begin{note}[title=Lendo o código de cores]
Em todo o software, a cor indica o nível de confiança associado a um número.
\textbf{Verde} = validado ou objetivo. \textbf{Âmbar} = experimental ou uma
estimativa, a ser validada depois. \textbf{Vermelho} é reservado para pareamento
quebrado e nunca usado para uma série de dados. Todo valor colorido carrega um
rótulo: a cor nunca é o único sinal.
\end{note}
